problem:  
Let \((X^n,d_X)\) be a compact Alexandrov space with curvature \(\ge 0\).  
Suppose there exists a *submetry*  
\[
\pi : X \to Y
\]
onto another compact Alexandrov space \((Y^{n-k}, d_Y)\) with curvature \(\ge 0\), such that each fiber \(F_y = \pi^{-1}(y)\) is a compact \(k\)-dimensional **totally geodesic Alexandrov subspace** (meaning every geodesic segment in \(F_y\) remains locally minimizing in \(X\)).  
Assume additionally that the regular parts \(M_{\mathrm{reg}}(X)\) and \(M_{\mathrm{reg}}(Y)\) are smooth Riemannian manifolds with nonnegative sectional curvature and that the singular sets of both spaces have Hausdorff codimension at least \(3\).  

**Question:**  
Does there exist a sequence of smooth Riemannian submersions
\[
\pi_i : (M_i^n, g_i) \longrightarrow (N_i^{n-k}, h_i)
\]
between compact Riemannian manifolds with \(\sec(g_i)\ge 0\), \(\sec(h_i)\ge 0\), such that  
\[
(M_i,g_i) \xrightarrow{\mathrm{GH}} (X,d_X), \qquad (N_i,h_i) \xrightarrow{\mathrm{GH}} (Y,d_Y),
\]
and \(\pi_i\) converges (in the Gromov–Hausdorff sense of graphs) to the given submetry \(\pi\)?  

Equivalently: is every compact submetry between nonnegatively curved Alexandrov spaces with totally geodesic fibers realizable as the limit of smooth nonnegatively curved Riemannian submersions?

---

Why is it a "good" problem:  

1. **Conceptual depth:**  
   This question directly probes the frontier of how metric curvature bounds interact with smooth geometric structure. Alexandrov submetries generalize Riemannian submersions, but it is unknown whether they always arise as geometric limits of smooth submersions under nonnegative curvature. Resolving this would effectively characterize when purely metric submersion structures encode smooth Riemannian origins—an outstanding issue in the transition from metric to smooth geometry.

2. **Bridge between fields:**  
   The problem connects several major areas of contemporary geometry:
   - **Comparison geometry**: via nonnegative Alexandrov curvature and totally geodesic subsets;  
   - **Riemannian submersion theory**: through O’Neill’s curvature formulas and fiber geometry;  
   - **Metric limit theory**: through Gromov–Hausdorff convergence and Perelman–Yamaguchi–Cheeger–Fukaya structural results.  
   Its resolution would require and unify tools from all three domains.

3. **Potential significance:**  
   - A *positive resolution* would create a powerful “metric-to-smooth lifting principle,” showing that every finely structured submetry in Alexandrov geometry comes from an underlying smooth nonnegatively curved model.  
   - A *counterexample* would identify genuinely new metric phenomena—“submetry fibrations” with no smooth analogues—thus revealing intrinsic geometric features that survive curvature comparison but defy smooth approximation.  
   Either outcome would substantially refine our understanding of curvature, topology, and smoothability.

4. **Simplicity and profundity:**  
   The statement is concise and purely geometric: “Does every nonnegatively curved Alexandrov submetry come from a smooth nonnegatively curved submersion in the GH limit?” Yet its resolution would require deep insight into curvature comparison, limit stability, and the fine geometry of singular spaces. The problem’s minimal phrasing but far-reaching implications exemplify the “narrow entrance, profound depth” aesthetic of truly great mathematical questions.